% -*- TeX-master: "../main.tex" -*-
\chapter{Introducción}

Los lenguajes de programación probabilística son lenguajes de
programación que incluyen alguna operación o construcción incorporada
que facilita la implementación de programas probabilísticos. Estos son
programas que representan una distribución de probabilidad implícita y
que, en última instancia, pueden reducirse a una representación
explícita de dicha distribución o a una muestra tomada de ella
\citep{prob-prog}. Los lenguajes de programación probabilística han
adquirido una importancia creciente debido a sus aplicaciones en
inferencia bayesiana y aprendizaje automático.

La mayoría de los lenguajes de programación poseen alguna forma de
tipado que, en general, cumple funciones como la detección de errores,
la restricción en la interacción entre módulos, la documentación de
funciones y variables, el aumento de la seguridad del lenguaje y la
facilitación del razonamiento sobre los programas
\citep{pierce}. Dependiendo de si la verificación de tipos se realiza
durante el tiempo de compilación o en tiempo de ejecución, se habla de
tipado estático o dinámico, respectivamente. Cada sistema tiene sus
ventajas y desventajas: el tipado estático permite detectar errores de
manera temprana, pero puede resultar demasiado restrictivo durante el
desarrollo de prototipos, mientras que el tipado dinámico permite un
desarrollo más rápido, a costa de un tiempo de ejecución más lento y
de la posibilidad de que algunos errores pasen desapercibidos.

El tipado gradual~\citep{gt} ofrece un equilibrio entre las dos
formas de tipado, mediante la introducción de tipos parcialmente
conocidos, que pueden ir desde un tipo completamente desconocido
($\Any$), hasta un tipo completamente preciso (por ejemplo, $\Int$),
incluyendo cualquier tipo compuesto parcialmente conocido (por
ejemplo, $\Fun{\Int}{\Any}$). En tiempo de compilación, el sistema de
tipos realiza un análisis estático para detectar posibles
inconsistencias de tipos. Si el programa es consistente, se asume que
está bien tipado, y el sistema comprueba en tiempo de ejecución si
esta suposición se mantiene. De este modo, puede lograrse cualquier
proporción de tipado estático y gradual (incluyendo tipado totalmente
estático o totalmente dinámico).

Se han desarrollado variantes gradualmente tipadas para lenguajes con
una amplia variedad de sistemas de tipos: orientados a objetos
\citep{siek2007,taki2012}, con efectos \citep{banados2014},
con tipado de seguridad \citep{fennell2013}, con control de seguridad
\citep{toro2018}, con referencias mutables \citep{gt,siek2015c,toro2020}, con tipos algebraicos de datos
\citep{malewski2021}, entre otros. Sin embargo, hasta hace poco no se
había intentado desarrollar un lenguaje probabilístico gradual.

El trabajo de \cite{gplc} introdujo el primer lenguaje probabilístico
gradual, el cálculo lambda probabilístico gradual o GPLC (por sus
siglas en inglés, \emph{Gradual Probabilistic Lambda Calculus}), que
en su extremo estático corresponde a un cálculo lambda equipado con
una operación de elección probabilística binaria. Este lenguaje
estático se extiende al lenguaje gradual GPLC mediante la metodología
AGT, propuesta por \cite{agt}.


La relevancia de GPLC no radica únicamente en proponer una nueva
combinación de ideas ya conocidas, sino en reunir en un mismo lenguaje
dos dimensiones de incertidumbre que suelen estudiarse por separado:
la incertidumbre propia del cómputo probabilístico y la incertidumbre
asociada a la información de tipos parcial o no especificada. Esta
combinación vuelve especialmente interesante estudiar no solo la
formulación formal del lenguaje, sino también su comportamiento
efectivo al ejecutar programas concretos, observar sus resultados y
contrastarlos con la intuición que motiva su diseño.

En este contexto, la construcción de un intérprete cumple un papel que
va más allá de la mera implementación de software. Una formulación
teórica permite definir con precisión la sintaxis, el sistema de tipos
y las reglas de evaluación de un lenguaje, pero una herramienta
ejecutable obliga a traducir esas definiciones a estructuras de datos,
algoritmos y mecanismos concretos de procesamiento. Esa traducción
hace posible someter el lenguaje a experimentación sistemática,
contrastar casos de uso reales, detectar ambigüedades o tensiones de
diseño y evaluar de manera más directa el alcance práctico de la
propuesta formal.

La presente memoria aborda ese problema mediante el desarrollo de un
intérprete prototipo de GPLC implementado en
OCaml~\citep{ocaml-system}. La elección de este lenguaje responde a
que ofrece un entorno especialmente adecuado para la implementación de
lenguajes de programación, tanto por su soporte para tipos algebraicos
y \emph{pattern matching} como por su utilidad para representar y manipular
estructuras recursivas como el AST.

El objetivo principal de esta memoria es desarrollar un intérprete
prototipo para el lenguaje GPLC, que implemente fielmente la semántica
de distribución descrita por \cite{gplc} y permita la ejecución
práctica de programas probabilísticos graduales. Con este intérprete
se realizará una validación empírica de la semántica existente
mediante la ejecución de ejemplos con diversas combinaciones de
constructos y distintos niveles de anotación gradual, de modo que se
compruebe su corrección y se observe su comportamiento en escenarios
reales.

De este modo, el aporte de esta memoria consiste en implementar la
propuesta formal de GPLC en un intérprete. En particular, esta memoria
busca aportar una herramienta que permita ejecutar programas que antes
existían únicamente como objetos de estudio teórico, y que a la vez
sirva como plataforma de experimentación para analizar el
comportamiento del lenguaje, examinar los alcances de su sistema de
tipos y producir evidencia concreta sobre la consistencia de su diseño
semántico.

Al mismo tiempo, la implementación permite evaluar de manera más
precisa qué aspectos de la formulación pueden implementarse de manera
directa en el intérprete y cuáles introducen dificultades no evidentes
en la exposición teórica. En este sentido, el intérprete no solo
constituye una herramienta para ejecutar programas, sino también un
medio para poner a prueba la teoría desde la práctica,
retroalimentando la comprensión del lenguaje a partir de su
comportamiento observable.

El resto de esta memoria se organiza de la siguiente manera. El
Capítulo 2 presenta el marco teórico necesario sobre lenguajes
probabilísticos, sistemas de tipos estáticos y tipado gradual. El
Capítulo 3 revisa el estado del arte más directamente relacionado con
implementaciones graduales y con el trabajo teórico en que se basa
GPLC. El Capítulo 4 expone el lenguaje GPLC desde su formulación
teórica. El Capítulo 5 presenta su sintaxis concreta y describe el
comportamiento general del intérprete desarrollado. El Capítulo 6
detalla la arquitectura del intérprete, mientras que el Capítulo 7
aborda las decisiones principales de implementación. El Capítulo 8
presenta la validación experimental mediante programas representativos
y discute las limitaciones observadas durante el
desarrollo. Finalmente, el Capítulo 9 expone las conclusiones del
trabajo y las principales líneas de trabajo futuro.

